% Stage 7 supplement: graph, geometric, and counting completion.
% Existing chapters remain authoritative; this module supplies bridges and
% student-facing algorithms that connect their techniques.
\chapter{Graph, Geometry, and Counting Completion}
\label{ch:stage7-graph-geometry-counting}

This chapter completes a practical bridge between randomized graph algorithms,
computational geometry, and approximate counting.  The recurring student
contract is explicit: identify the random experiment, state what it estimates
or preserves, quantify failure, and account for the cost of repetition.  The
earlier chapters on min-cut, geometry, Markov chains, and approximate counting
remain in place; the present chapter connects them and adds several standard
techniques that are easy to use incorrectly in isolation.

\section{Recursive min-cut and network reliability}

Karger's contraction algorithm preserves a fixed minimum cut with probability
roughly $\Theta(n^{-2})$.  Repetition therefore gives a simple high-probability
algorithm, but recursive contraction gives a better way to spend trials.  Keep
contracting a random edge until $r=\lceil n/\sqrt{2}\rceil+1$ supervertices
remain, recursively solve two independent copies, and return the better cut.
Conditioned on a fixed minimum cut surviving, each recursive child has a
constant chance of finding it.  The resulting recurrence is
\[
 T(n) \leq 2T(\lceil n/\sqrt{2}\rceil+1)+O(n^2),
\]
which yields $T(n)=O(n^2\log n)$ for a straightforward implementation (and
the usual $O(n^2\log^3 n)$ high-probability bound after amplification).

The analysis teaches an important distinction: the algorithm may return a cut
that is valid but not minimum.  Its one-sided failure event is ``the target
cut was destroyed or never recovered.''  A student implementation should
test the returned cut's value, use multigraph edges with multiplicity, and
never sample uniformly from distinct endpoint pairs.

For a graph with independent edge-failure probabilities, network reliability
is the probability that the surviving graph connects the designated terminals.
Exact evaluation is generally difficult, but Monte Carlo is immediate:
sample a failed/surviving edge set, run connectivity, and average the indicator.
If the reliability is $p$, then $N$ samples have variance $p(1-p)/N$ and
Hoeffding gives
\[
 \Pr[|\widehat p-p|>\varepsilon]\leq 2e^{-2N\varepsilon^2}.
\]
This is an additive guarantee.  When $p$ is very small, it may require far
too many samples; importance sampling or splitting is needed for relative
error.  This is the first rare-event warning in the chapter.

\section{Random spanning trees and Kirchhoff's theorem}

The uniform spanning tree (UST) of a connected graph is a tree selected
uniformly from all spanning trees.  Kirchhoff's matrix-tree theorem provides
the count:
\[
 \tau(G)=\det L^{(v)},
\]
where $L$ is the graph Laplacian and $L^{(v)}$ deletes row and column $v$.
The formula is useful even when it is not an algorithm: it connects counting,
determinants, electrical networks, and sampling.

Wilson's algorithm samples a UST exactly.  Start with a root in the tree;
from each unvisited vertex, run a simple random walk until it hits the tree,
erase loops chronologically from the walk, and add the resulting loop-erased
path.  The output is independent of the order in which vertices are started.
The walk can be long, so the expected running time is naturally expressed in
terms of hitting times; it is not automatically linear in the number of edges.

For an edge $e$, the UST inclusion probability satisfies
\[
 \Pr[e\in T]=w_e R_{\mathrm{eff}}(e),
\]
where $w_e$ is its conductance and $R_{\mathrm{eff}}(e)$ is effective
resistance between its endpoints.  Thus sampling USTs gives unbiased
estimators for network quantities, while Kirchhoff's theorem gives exact
answers on graphs small enough for determinant computation.

\section{Metric embeddings}

An embedding maps a finite metric $(X,d)$ into a simpler metric space while
controlling distortion.  A randomized embedding into a tree metric is useful
because many optimization problems are easier on trees.  Bourgain's embedding
uses, for each scale, random subsets $S$ and the coordinates
\[
 x_{S}(u)=d(u,S)=\min_{s\in S}d(u,s).
\]
After $O(\log |X|)$ independent scales and suitable normalization, the map
into $O(\log^2 |X|)$ dimensions has distortion $O(\log |X|)$.  The proof has
two parts: distance-to-a-set is 1-Lipschitz, giving non-contraction, and a
random subset separates any pair at an appropriate scale with enough
probability, giving the expansion bound.

The practical lesson is to state the direction of each inequality.  A map may
be non-contracting but have large expansion, or it may preserve distances in
expectation without preserving every pair simultaneously.  For an algorithm,
the embedding dimension, distortion, construction time, and probability that
all pairwise bounds hold are separate parameters.

\section{Randomized point location}

In a planar subdivision, point location asks which face contains a query point.
A randomized incremental trapezoidal map inserts segments in random order and
maintains the vertical decomposition.  The query structure is a directed
acyclic graph; a query follows the comparisons made during construction.
Backward analysis shows that the expected search depth is $O(\log n)$ when
the conflict history is handled correctly.  The same pattern appears in
randomized incremental convex hulls and Delaunay triangulations:
\begin{enumerate}
\item define the conflict set of a feature;
\item expose the last inserted object affecting the query;
\item bound the probability that this object is among the few relevant ones;
\item sum the expected update cost.
\end{enumerate}
Degeneracies are not a cosmetic detail.  A robust implementation needs a
consistent tie-breaking rule (often symbolic perturbation) for collinear
points, endpoint queries, and overlapping boundaries.  Without it, the
mathematical probability analysis may describe an algorithm whose comparison
tree is not even well-defined.

\section{Rare-event estimation}

Suppose an event $A$ has probability $p\ll1$.  The indicator estimator is
unbiased but has relative standard deviation
\[
 \frac{\sqrt{\operatorname{Var}(1_A)}}{\mathbb E[1_A]}
 =\sqrt{\frac{1-p}{p}},
\]
so estimating $p$ to fixed relative precision needs $\Theta(1/p)$ samples.
Importance sampling changes the sampling distribution from $P$ to $Q$ and
uses the unbiased weight
\[
 \widehat p=\frac1N\sum_{i=1}^N
 1_A(X_i)\frac{P(X_i)}{Q(X_i)},\qquad X_i\sim Q.
\]
The support condition $Q(x)>0$ whenever $P(x)1_A(x)>0$ is mandatory.  A poor
choice of $Q$ can have infinite or enormous variance, so the student should
report an empirical variance estimate and compare it with crude Monte Carlo.
Splitting and conditional Monte Carlo are alternative ways to replace a rare
indicator by a less variable conditional expectation.

\section{Counting versus generation}

Let $\Omega$ be a finite family and suppose a sampler produces an exact
uniform element of $\Omega$.  Sampling alone does not reveal $|\Omega|$;
generation and counting are different tasks.  The standard bridge is a chain
of nested sets
\[
 \Omega_0\supseteq\Omega_1\supseteq\cdots\supseteq\Omega_k,
 \qquad |\Omega_0|\text{ known},
\]
whose successive ratios are bounded away from zero.  If a sampler can estimate
each ratio $|\Omega_{i+1}|/|\Omega_i|$ with relative error and confidence,
then multiplying the estimates yields an FPRAS after allocating the error
budget across levels.  The product structure explains why warm starts,
overlap between adjacent distributions, and a polynomial number of levels are
essential.

Conversely, an approximate counter can often support generation by recursively
choosing a branch with probability proportional to its estimated number of
completions.  Approximation errors accumulate, so one must specify whether
the output is exactly uniform, approximately uniform in total variation, or
only correct in expectation.  These guarantees are not interchangeable.

\section{Student checklist and exercises}

For every graph, geometry, or counting algorithm, write down: the object being
sampled; the target distribution or estimator; the invariant; the failure
probability; the runtime and space; and the condition under which the claim is
valid.  In particular, separate an additive error claim from a relative one.

\begin{problem}[Reliability]
For independent edge failures, derive the variance of the connectivity
indicator and choose $N$ for additive error $\varepsilon$ and failure at most
$\delta$.  Explain why the same $N$ is not a relative-error guarantee when
the reliability is unknown and may be exponentially small.
\end{problem}

\begin{problem}[UST experiment]
Implement Wilson's algorithm on a small connected graph.  Compare empirical
edge inclusion frequencies with $w_eR_{\mathrm{eff}}(e)$ computed from a
Laplacian solve.
\end{problem}

\begin{problem}[Embedding audit]
For a metric on $n$ points, implement random distance-to-subset coordinates.
Measure the largest contraction and expansion over all pairs, and explain
which observed failures are due to finite sampling rather than the theorem's
asymptotic constants.
\end{problem}

\begin{problem}[Counting bridge]
Assume $k$ nested ratios each lie in $[1/4,1]$.  Allocate a total relative
error $\varepsilon$ and failure probability $\delta$ across the ratios and
derive a polynomial sample bound for their product.
\end{problem}
